\begin{trapbox}
	\begin{description}[style=nextline, leftmargin=2.8em, labelsep=0.5em, itemsep=0.35em]
		\item[\textbf{\textcolor{CTrap}{易错点一（忽略定义域）}}] 求对数型函数的定义域或解对数不等式时，忘记检查真数大于$0$。
		\item[\textbf{\textcolor{CTrap}{正确理解（先定义域）}}] 涉及对数函数的任何问题，必须先保证自变量满足真数$>0$。
		\item[\textbf{\textcolor{CTrap}{错因分析（习惯性疏忽）}}] 习惯将所有函数定义域默认为$\mathbb{R}$，遗忘对数函数的限制条件。
	\end{description}

	\medskip
	\begin{description}[style=nextline, leftmargin=2.8em, labelsep=0.5em, itemsep=0.35em]
		\item[\textbf{\textcolor{CTrap}{易错点二（底数分类遗漏）}}] 处理含有参数$a$的对数函数单调性或不等式时，直接假设底数$>1$，未分$0<a<1$讨论。
		\item[\textbf{\textcolor{CTrap}{正确理解（分类讨论）}}] 当底数含参数$a$时，必须分$a>1$和$0<a<1$两类分别处理。
		\item[\textbf{\textcolor{CTrap}{错因分析（默认递增）}}] 默认对数函数都是增函数，忽略底数在$(0,1)$时的递减情形。
	\end{description}

	\medskip
	\begin{description}[style=nextline, leftmargin=2.8em, labelsep=0.5em, itemsep=0.35em]
		\item[\textbf{\textcolor{CTrap}{易错点三（符号记忆混淆）}}] 记错函数值符号，比如误以为$a>1$时$x>1$对应$y<0$。
		\item[\textbf{\textcolor{CTrap}{正确理解（符号法则）}}] $a>1$时，$x>1\Rightarrow y>0$，$0<x<1\Rightarrow y<0$；$0<a<1$时反之。
		\item[\textbf{\textcolor{CTrap}{错因分析（机械记忆）}}] 不结合单调性与定点推导符号，死记硬背导致混淆。
	\end{description}
\end{trapbox}
